\section{Extracting a derivation on the determinant ring}
\label{sec:reduction}

We now pass from a homogeneous additive action on the degeneration to a
locally nilpotent derivation of
\[
  \Rt=\kk[a,d,b,c,u]/(ac-bd-1)
\]
with a prescribed sign for its fiber weight.  Throughout this section,
\emph{degree} means the valuation degree, so that
\[
  \deg \Rt=0,\qquad \deg\tau=-1,\qquad \deg V=2.
\]
The separate fiber weights are
\[
  \wt(a)=\wt(d)=1,\quad \wt(b)=\wt(c)=-1,\quad
  \wt(u)=\wt(\tau)=0,\quad \wt(V)=2.
\]
By \Cref{prop:principalization}, we work in the domain
\begin{equation}\label{red:ring}
  \Gt=\Rt[\tau,V]\subset\Rt[\tau,\tau^{-1}],
  \qquad \tau^2V=v=a^3b-a^2u^3-d^2.
\end{equation}
In particular, $\tau$, $V$, and $v$ are nonzero.  The notation
$\Rt[\tau,V]$ here denotes the indicated subalgebra, with the displayed
relation.

The positive-degree invariant supplied by polynomiality will be
divisible by $V$. The elementary order of a locally nilpotent
derivation makes its kernel factorially closed, so this divisibility
will force $V$ itself to be invariant. We recall the argument,
including the order estimates needed when $\tau$ moves; compare
Finston--Maubach~\cite[Lemma~1.1]{FinstonMaubach}.

\begin{lemma}[Order of a locally nilpotent derivation]
\label{lem:lnd-order}
Let $T$ be a domain over $\kk$ and let $\delta$ be a locally nilpotent
derivation of $T$.  For $0\ne h\in T$, define
\[
  \nu_\delta(h)=\max\{n\ge0:\delta^n(h)\ne0\}.
\]
Then for nonzero $f,g\in T$,
\begin{equation}\label{red:order-product}
  \nu_\delta(fg)=\nu_\delta(f)+\nu_\delta(g).
\end{equation}
If $\delta(h)\ne0$, then
$\nu_\delta(\delta(h))=\nu_\delta(h)-1$.
Consequently, $\ker\delta$ is factorially closed in $T$: if
$0\ne fg\in\ker\delta$, then $f,g\in\ker\delta$.
Moreover, $\delta(h)\ne0$ implies $\delta(h)\notin hT$.
\end{lemma}

\begin{proof}
Local nilpotence makes the defining maximum finite.  Set
$r=\nu_\delta(f)$ and $s=\nu_\delta(g)$.  The iterated Leibniz rule gives
\[
  \delta^n(fg)=\sum_{i=0}^n\binom ni\delta^i(f)\delta^{n-i}(g).
\]
For $n>r+s$ every term is zero.  For $n=r+s$ only the term with $i=r$
can be nonzero, and it equals
\[
  \binom{r+s}{r}\delta^r(f)\delta^s(g)\ne0.
\]
Here we use both characteristic zero and the domain property.  This
proves \Cref{red:order-product}.  The assertion about $\delta(h)$ follows
directly from the definition.  Since a nonzero element has order zero
exactly when it is invariant, product additivity proves factorial
closure.  Finally, if $0\ne\delta(h)=hq$, then $q\ne0$ and
\[
  \nu_\delta(h)-1
    =\nu_\delta(\delta(h))
    =\nu_\delta(h)+\nu_\delta(q)
    \ge\nu_\delta(h),
\]
which is impossible.
\end{proof}

We will also use the exact homogeneous pieces of \Cref{red:ring}.
Writing $\Gt_e$ for its degree-$e$ component, they are
\begin{equation}\label{red:pieces}
  \Gt_e=
  \begin{cases}
    \tau^{-e}\Rt, & e<0,\\
    V^\ell\Rt, & e=2\ell,\quad \ell\ge0,\\
    \tau V^{\ell+1}\Rt, & e=2\ell+1,\quad \ell\ge0.
  \end{cases}
\end{equation}
Indeed, a degree-$e$ element is an $\Rt$-linear combination of monomials
$\tau^iV^h$ with $i,h\ge0$ and $2h-i=e$.  Using $\tau^2V=v$,
these monomials become, respectively,
\[
  \tau^{-e}v^h,\qquad
  V^\ell v^{h-\ell},\qquad
  \tau V^{\ell+1}v^{h-\ell-1}.
\]
The exponents of $v$ in the relevant cases are nonnegative.  Conversely,
the multipliers displayed in \Cref{red:pieces} belong to $\Gt$ and have
the required degrees.  This proves the equalities, including that their
coefficients lie in $\Rt$.  Since the multipliers are nonzero and $\Gt$
is a domain, each such coefficient is unique.

\begin{proposition}\label{prop:extracted-lnd}
If $A\cong\kk\poly{4}$, there are a nonzero locally nilpotent derivation
$E$ of $\Rt$ and an integer $q\le0$ such that $E$ is homogeneous of
fiber weight shift $q$ and
\[
  E^2(v)=0.
\]
Here homogeneity means that $E(r)$ has fiber weight $\wt(r)+q$ whenever
$r$ is fiber homogeneous and $E(r)\ne0$.
\end{proposition}

\begin{proof}
By \Cref{prop:positive-invariant}, polynomiality of $A$ supplies a nonzero
homogeneous LND $D_0$ of $G$ and a nonzero homogeneous
$h_0\in\ker D_0$ of positive degree.  Apply
\Cref{prop:homogeneous-lift} to obtain an LND $\Dt$ of $\Gt$ lifting
$D_0$, homogeneous of some degree shift $j\in\mathbb Z$ and of fiber
weight shift zero.  The injection $G\hookrightarrow\Gt$ in
\Cref{prop:principalization} ensures both that $\Dt\ne0$ and that the
image $h$ of $h_0$ is nonzero.  This image is a homogeneous invariant
of the same positive degree.  By \Cref{red:pieces}, $h$ is divisible
by $V$ in $\Gt$.  Factorial closure in \Cref{lem:lnd-order} therefore
gives
\begin{equation}\label{red:V-invariant}
  \Dt(V)=0.
\end{equation}
We distinguish whether $\tau$ is also invariant. If it is, the
derivation restricts after removing a power of $\tau$. If it is not,
the homogeneous pieces force a positive factor of $V$ into its values
on $\Rt$; removing this factor produces a derivation of negative
fiber weight. The LND order proves local nilpotence in this second
case without assuming that $\Dt$ preserves $\Rt$.

\medskip
\noindent\textbf{Case 1: $\Dt(\tau)=0$.}
In this case $\Dt$ extends to a locally nilpotent derivation of
\[
  \Gt[\tau^{-1}]=\Rt[\tau,\tau^{-1}].
\]
Indeed, for $f\in\Gt$ and $m\ge0$, all iterates on
$f\tau^{-m}$ equal $\tau^{-m}\Dt^n(f)$.  The degree-$j$ component of
this Laurent polynomial ring is exactly $\tau^{-j}\Rt$.
Since every $r\in\Rt$ has degree zero, there is thus a unique
$E(r)\in\Rt$ satisfying
\begin{equation}\label{red:fixed-tau}
  \Dt(r)=\tau^{-j}E(r).
\end{equation}
Linearity and the Leibniz rule for $\Dt$, followed by cancellation of
$\tau^{-j}$, show that $E:\Rt\to\Rt$ is a $\kk$-derivation.
It has fiber weight shift zero because $\tau$ has fiber weight zero
and $\Dt$ preserves fiber weight.  Since $\Dt(\tau)=0$, induction in
\Cref{red:fixed-tau} gives
\[
  \Dt^n(r)=\tau^{-nj}E^n(r)\qquad(n\ge0).
\]
Local nilpotence of $\Dt$ therefore proves local nilpotence of $E$.
By \Cref{red:V-invariant},
\[
  0=\Dt(\tau^2V)=\Dt(v)=\tau^{-j}E(v),
\]
so $E(v)=0$.  Finally, $E=0$ would make $\Dt$ vanish on $\Rt$, $\tau$,
and $V$, which generate $\Gt$; hence $E\ne0$.  This proves the desired
conclusion with $q=0$.

\medskip
\noindent\textbf{Case 2: $\Dt(\tau)\ne0$.}
Write $\nu=\nu_{\Dt}$ on the nonzero elements of $\Gt$, and put
$m=\nu(\tau)\ge1$.  By \Cref{lem:lnd-order},
$\Dt(\tau)\notin\tau\Gt$.  On the other hand, this element is
homogeneous of degree $j-1$.  \Cref{red:pieces} shows that
every component of negative degree, and every component of positive
odd degree, is contained in $\tau\Gt$.  We must therefore have
\begin{equation}\label{red:odd-shift}
  j-1=2\ell\quad\text{for some }\ell\ge0,
  \qquad j=2\ell+1.
\end{equation}
For $r\in\Rt$, the value $\Dt(r)$ belongs to the degree-$j$ piece.
By \Cref{red:pieces} there is a unique element $E(r)\in\Rt$ such that
\begin{equation}\label{red:moving-tau}
  \Dt(r)=\tau V^{\ell+1}E(r).
\end{equation}
Thus the values of $E$ lie in $\Rt$ before any iteration or nilpotence
argument is used.  Cancellation of the common nonzero multiplier in
the linearity and Leibniz identities again proves that $E$ is a
$\kk$-derivation of $\Rt$.  Since that multiplier has fiber weight
$2(\ell+1)$, the fiber weight shift of $E$ is
\begin{equation}\label{red:negative-weight}
  q=-2(\ell+1)<0.
\end{equation}

We prove local nilpotence inside $\Gt$ using its LND order.  Whenever
$E(r)\ne0$, \Cref{red:moving-tau}, \Cref{red:V-invariant}, and
\Cref{lem:lnd-order} imply
\begin{equation}\label{red:order-decrease}
  \nu(E(r))=\nu(r)-1-m.
\end{equation}
In fact, the order of the left side of \Cref{red:moving-tau} is
$\nu(r)-1$, whereas that of its right side is
$m+(\ell+1)\nu(V)+\nu(E(r))=m+\nu(E(r))$.
Every iterate $E^n(r)$ still lies in $\Rt$.  Consequently, as long as
it is nonzero, successive applications of \Cref{red:order-decrease}
give
\[
  \nu(E^n(r))=\nu(r)-n(m+1).
\]
This nonnegative integer cannot decrease indefinitely, proving that
$E$ is locally nilpotent.

Differentiating the relation in \Cref{red:ring} gives
\begin{equation}\label{red:moving-v}
  \tau V^{\ell+1}E(v)=\Dt(v)=2\tau V\Dt(\tau)\ne0.
\end{equation}
The last inequality uses characteristic zero and the domain property.
Thus $E(v)\ne0$, and in particular $E\ne0$.
Product additivity and \Cref{red:V-invariant} give
$\nu(v)=\nu(\tau^2V)=2m$; hence \Cref{red:order-decrease} yields
\[
  \nu(E(v))=m-1.
\]
If $E^2(v)$ were nonzero, a second application would yield
\[
  \nu(E^2(v))=(m-1)-1-m=-2,
\]
an impossibility.  Therefore $E^2(v)=0$, as required.
\end{proof}

\begin{remark}\label{red:sign-essential}
The restriction to nonpositive fiber weight is essential to the
obstruction proved in \Cref{sec:rigidity}.  Indeed, the derivation
\[
  E_+(b)=a,\qquad E_+(c)=d,\qquad
  E_+(a)=E_+(d)=E_+(u)=0
\]
annihilates $ac-bd-1$ and so defines a derivation of $\Rt$.
It is locally nilpotent, since on the polynomial ring it lowers total
degree in $b,c$, and it has fiber weight shift $+2$.  Nevertheless,
$E_+(v)=a^4$ and $E_+^2(v)=0$.
\end{remark}
